\begin{figure*}[t]
\centering
\resizebox{\textwidth}{!}{%
\begin{tikzpicture}[
  font=\footnotesize,
  stage/.style={draw=black!40, rounded corners=2pt, fill=#1, align=left, inner sep=4pt, minimum height=3.1cm, anchor=north west},
  arr/.style={-{Stealth[length=2mm]}, thick, black!55},
]
\node[stage=blue!4, text width=3.6cm] (data) at (0,0)
  {\textbf{PaHaW recordings}\\[2pt]
   75 participants, 8 tasks,\\597 recordings\\[3pt]
   per sample: position, time, pen contact, pressure, altitude and azimuth};

\node[stage=orange!7, text width=4.0cm, right=5mm of data.north east, anchor=north west] (render)
  {\textbf{Pseudodynamic minimap}\\[2pt]
   \textbullet\ trace drawn only where the pen touches the tablet\\
   \textbullet\ each signal normalized within its own recording\\
   \textbullet\ azimuth centered on its circular mean\\
   \textbullet\ up to three signals mapped to RGB};

\node[stage=orange!7, text width=3.9cm, right=5mm of render.north east, anchor=north west] (enc)
  {\textbf{Nine encodings} (Table~\ref{tab:encoding})\\[3pt]
   \begin{tabular}{@{}c@{\hspace{8pt}}c@{}}
   \includegraphics[height=1.2cm]{method_static.png} & \includegraphics[height=1.2cm]{method_saz.png}\\
   \scriptsize 1: static trace & \scriptsize 8: $SAZ$
   \end{tabular}\\[2pt]
   2--5: one signal; 6--9: RGB triplets};

\node[stage=blue!4, text width=3.7cm, right=5mm of enc.north east, anchor=north west] (eval)
  {\textbf{Evaluation}\\[2pt]
   $5\times5$ participant-separated cross-validation with 3 inner folds\\[3pt]
   recording-level macro F1\\[1pt]
   participant-level macro F1 (majority vote over tasks)};

\draw[arr] ([yshift=-1.55cm]data.north east) -- ([yshift=-1.55cm]render.north west);
\draw[arr] ([yshift=-1.55cm]render.north east) -- ([yshift=-1.55cm]enc.north west);
\draw[arr] ([yshift=-1.55cm]enc.north east) -- ([yshift=-1.55cm]eval.north west);
\end{tikzpicture}}
\caption{Overview of the method. Each PaHaW recording is rendered as a pseudodynamic minimap in one of nine encodings, and the three experiments share the same participant-separated cross-validation.}
\label{fig:method}
\end{figure*}
